%% toc: P3 26 26.1,26.2,26.3
\chapter{Linear Orders as the Model Case of Fra\"iss\'e's Relation Theory}
\label{sec:p03-w03-p03-014-orders}

Fra\"iss\'e introduced his classes of relations as a generalization of
linear orders and of finite linear orders, and he described several of his
general theorems as extensions of familiar facts about orders
\cite[\S1.1, pp.~363--364]{Fraisse1954}. This chapter makes that
specialization explicit. It derives each order fact from the general
results of \cref{sec:p03-fraisse-1954}, checks every hypothesis used, and
records the order-theoretic remarks scattered through the source:
the local class of an infinite order, orders on arbitrary sets,
homogeneous orders and the rational order, and the witnesses of the
countable universality condition.

\emph{Inputs.} A linear order, presented as in the source by its reflexive
two-valued relation (\cref{def:p03-fraisse54-relation}); the closure
conditions, local classes and homogeneity notions of
\cref{def:p03-fraisse54-age,def:p03-fraisse54-closures,def:p03-fraisse54-extension};
and the general theorems \cref{thm:p03-fraisse54-gamma,thm:p03-fraisse54-extension,thm:p03-fraisse54-original-construction,thm:p03-fraisse54-orders}.
\emph{Conclusions.} The local class of an infinite order is the class of
all linear orders; every set carries a linear order; two orders, or any
set of orders, embed into one order; every countable order embeds into
the rational order; the rational order is the unique countable homogeneous
order with more than one element; every infinite homogeneous order contains
a copy of the rational order; and the countable universal hosts for
countable orders are exactly the countable orders containing such a copy.
The ambient set theory is ZFC; Choice enters through the ultrafilters of
\cref{thm:p03-fraisse54-gamma}.

\section{The local class of an infinite order}
\label{sec:p03-w03-p03-014-local-class}

\begin{hbconvention}{def:p03-w03-p03-014-eta}{Orders, order sums, and \(\eta\)}
A linear order on \(E\) is a reflexive, antisymmetric, transitive and
total relation \(R:E^2\to\{+,-\}\), where \(R(x,y)=+\) means
\(x\leq y\). An embedding of orders is an embedding of these relations
in the sense of \cref{def:p03-fraisse54-relation}; for linear orders it is
the same as a strictly increasing map.
Following the source, \(\eta\) denotes the order type of the rational
numbers \((\mathbb Q,\leq)\), and ``\(R\) extends \(\eta\)'' means
that \(\eta\ll R\), that is, some copy of \(\eta\) is a restriction of
\(R\) \cite[\S\S1.1, 5.1, 6.3]{Fraisse1954}.
For orders \(A,B\) on disjoint bases, the sum \(A+B\) is the order on the
union that agrees with \(A\) and with \(B\) on their bases and puts every
element of the base of \(A\) below every element of the base of \(B\).
\end{hbconvention}

\begin{hbtheorem}{thm:p03-w03-p03-014-chain-age}{proved=annex:p03-w03-p03-014-chain-age}{The local class of an infinite order}
Let \(R\) be an infinite linear order. Then the age \(\gamma_R\) is
the class of all finite linear orders, and \(\Gamma_R\) is the class of
all linear orders.
\end{hbtheorem}

This is the source's example in \S1.4, recalled again in \S5.1
\cite[\S1.4, p.~365; \S5.1, p.~377]{Fraisse1954}.
For a finite order with \(p\) elements, by contrast, \(\Gamma_R\)
consists of the linear orders with at most \(p\) elements.

\begin{hbcorollary}{thm:p03-w03-p03-014-ordering}{proved=annex:p03-w03-p03-014-ordering}{Every set can be linearly ordered}
In set theory with Choice, every set carries a linear order.
This is obtained from the size alternative of
\cref{thm:p03-fraisse54-gamma} applied to \(\Gamma_R\) for an infinite
order \(R\).
\end{hbcorollary}

The source makes this remark immediately after its proof of Theorem~III,
part~1, and announces it in its opening as the fact that Theorem~III,
part~1, generalizes \cite[\S1.1, p.~363; \S4.1, p.~374]{Fraisse1954}.
The derivation is not a new proof of the ordering principle from weaker
assumptions: the normal ultrafilter on maps into an infinite base, which
the size alternative uses, is itself obtained with Choice.

\section{The four order propositions behind the general theorems}
\label{sec:p03-w03-p03-014-four}

\begin{hbtheorem}{thm:p03-w03-p03-014-four-facts}{proved=annex:p03-w03-p03-014-four-facts}{Order instances of Theorems III, IV, V and VI}
\begin{enumerate}
\item For linear orders \(A,B\), there is a linear order \(C\) with
      \(A\ll C\) and \(B\ll C\); after renaming the bases apart,
      \(C=A+B\) is one such order.
\item For every set of linear orders there is one linear order into which
      each of them embeds.
\item Every finite or countable linear order embeds into \(\eta\).
\item Up to isomorphism, \(\eta\) is the only nonempty countable linear
      order that is dense and has neither a first nor a last element. Among
      countable linear orders with more than one element, it is the only
      homogeneous one.
\end{enumerate}
Part~1 is the joint-embedding clause \(\mathrm C_2\) of Theorem~III, part~2,
for \(\Gamma_R\); part~2 is Theorem~IV; part~3 is Theorem~V; part~4 is
the uniqueness clause of Theorem~VI, part~2, together with the order
criterion of \S6.2. These are the correspondences stated in
\cite[\S1.1, pp.~363--364; \S5.1, p.~377; \S8.1, p.~383]{Fraisse1954}.
\end{hbtheorem}

In this chapter Theorems~III and IV are
\cref{thm:p03-fraisse54-gamma}, Theorem~V is parts~1 and 2 of
\cref{thm:p03-fraisse54-extension}, and Theorem~VI is
\cref{thm:p03-fraisse54-original-construction}.
Part~2 of \cref{thm:p03-w03-p03-014-four-facts} gives no bound on the size
of the common host; for countably many countable orders, part~3 provides
the countable host \(\eta\).

\section{Homogeneous orders and the rational order}
\label{sec:p03-w03-p03-014-homogeneous}

The source proves the facts of this section in two ways. The first route
passes through the extension condition \(\mathcal H\), using \S\S5.3, 5.5
and 6.2. The source adds that all of them can also be obtained directly
from the gap criterion of \S6.1 \cite[\S6.3, p.~380]{Fraisse1954}.
\Cref{thm:p03-w03-p03-014-direct-gap} supplies that direct route; its
proof does not use \(\mathcal H\).

\begin{hbtheorem}{thm:p03-w03-p03-014-direct-gap}{proved=annex:p03-w03-p03-014-direct-gap}{The direct route from the gap criterion}
Using only the gap criterion of \cref{thm:p03-fraisse54-orders} and the
countable uniqueness of \(\eta\):
\begin{enumerate}
\item a homogeneous linear order with more than one element is dense and
      has neither a first nor a last element;
\item a countable dense linear order with more than one element and
      without endpoints is homogeneous.
\end{enumerate}
\end{hbtheorem}

\begin{hbtheorem}{thm:p03-w03-p03-014-contains-eta}{proved=annex:p03-w03-p03-014-contains-eta}{Every infinite homogeneous order extends \(\eta\)}
Let \(I\) be an infinite homogeneous linear order. Then
\(\eta\ll I\). Moreover every finite or countable linear order embeds
into \(I\).
\end{hbtheorem}

This is the result announced in \S5.1 and drawn in \S6.3
\cite[\S5.1, p.~377; \S6.3, p.~380]{Fraisse1954}. It holds for
uncountable \(I\) as well, such as the real line or the uncountable
examples of \cref{thm:p03-fraisse54-orders}. The proof uses homogeneity
only through condition \(\mathcal H\), so the conclusion also holds for
an infinite order satisfying \(\mathcal H\) without being homogeneous,
such as the real line with one point removed.

\begin{hbtheorem}{thm:p03-w03-p03-014-condition-c}{proved=annex:p03-w03-p03-014-condition-c}{The witnesses of condition \(\mathcal C\) for orders}
The class of linear orders satisfies the countable universality condition
\(\mathcal C\) of \cref{def:p03-fraisse54-universal}. A finite or
countable linear order \(U\) witnesses \(\mathcal C\), that is, every
finite or countable linear order embeds into \(U\), if and only if
\(U\) is countable and \(\eta\ll U\).
\end{hbtheorem}

The ``if'' direction is the source's statement that \(\eta\) or any
countable order extending \(\eta\) serves as the host
\cite[\S4.7, p.~376]{Fraisse1954}; the converse shows that these are the
only hosts.

\begin{hbexample}{def:p03-w03-p03-014-example}{A non-homogeneous host and a non-host}
The countable order \(\omega+\eta\), a copy of the natural numbers
followed by a copy of the rationals, extends \(\eta\). By
\cref{thm:p03-w03-p03-014-condition-c} every countable order embeds
into it, so it witnesses \(\mathcal C\). It is not homogeneous: it has a
first element, which part~1 of \cref{thm:p03-w03-p03-014-direct-gap}
excludes. Thus a countable universal host need not be the homogeneous
representative of \cref{thm:p03-fraisse54-original-construction}; that
representative is singled out by homogeneity, not by universality.

The order \(\omega\) of the natural numbers is countable and infinite,
yet \(\eta\) does not embed into it: the image of \(\mathbb Q\) would be
a nonempty set of natural numbers without a least element. So \(\omega\)
is not a witness, and by \cref{thm:p03-w03-p03-014-contains-eta} it is
not homogeneous either. Directly, the interval between \(0\) and \(1\)
is empty while that between \(0\) and \(2\) is not, contrary to the
gap criterion.
\end{hbexample}

\begin{hbwarning}{def:p03-w03-p03-014-limits}{What the order case does not show}
The order facts above are instances, not proofs, of the general
theorems' strength. For linear orders the class has the strong
amalgamation condition \(\overline{\mathcal D}\)
(\cref{thm:p03-fraisse54-examples}), so it illustrates the construction in a class with this stronger
property. It does not exhibit the distinction between amalgamation and
strong amalgamation; the unary-relation example in
\cref{thm:p03-w03-p03-014-printed-unary} supplies that distinction. The ramified orders show the
other side: they form a local class without amalgamation, and the source
leaves open whether they satisfy \(\mathcal C\)
\cite[\S9.5, p.~387]{Fraisse1954}. The order instances say nothing about
that question.
\end{hbwarning}

\section{Relationships and sources}
\label{sec:p03-w03-p03-014-relations}

The general theory is \cref{sec:p03-fraisse-1954}; the modern
countable-signature form is \cref{thm:p03-fraisse54-countable-language},
and the age and ultrahomogeneity terminology is that of
\cref{def:p03-age-fraisse,def:p03-ultrahomogeneous}.
The first-order theory of dense linear orders without endpoints and its
countable categoricity are treated with the countable-categoricity
criterion in \cref{thm:p03-rn-dlo-uniqueness}. The source for every
order remark of this chapter is Fra\"iss\'e's 1954 article
\cite[\S\S1.1, 1.4, 4.1, 4.7, 5.1, 6.1--6.3, 8.1]{Fraisse1954}.
